% GENERATED FILE. Edit canonical lessons, then run npm run generate:books.
% canonical-source-hash: fnv1a64:ce6906dc414131e4
% canonical-lessons: FA-C501-baradarzade, FA-C501-madarshohar, FA-C501-pedarzan, FA-C501-doqolu, FA-C501-farzand

\chapter{In-laws and Twins}
\label{ch:fa-in-laws-and-twins}
\hlchaptermodality{501}

\begin{chapteropening}
\textbf{By the end of this chapter:} \emph{I can talk about in-laws and twins.}

Complete the last lesson of chapter 501: I can talk about in-laws and twins.
\end{chapteropening}

\section[\texorpdfstring{\fa{برادرزاده} (barâdarzâde)}{برادرزاده (barâdarzâde)}]{\fa{برادرزاده} (barâdarzâde) — a brother's child, a nephew or niece}

\label{lesson:FA-C501-baradarzade}



\begin{quote}
Before the new one: say the Persian for a maternal cousin, then the Persian for a sister's child.
\end{quote}

\subsection*{You'll want to know: \fa{برادرزاده}}
\textbf{\fa{برادرزاده}} (\emph{barâdarzâde}) — \textquotedblleft{}a brother's child, a nephew or niece\textquotedblright{}.

\begin{quote}
\fa{برادرزاده} \fa{من} \fa{تازه} \fa{راه} \fa{رفتن} \fa{یاد} \fa{گرفته} \fa{است}.
\end{quote}

\begin{quote}
\emph{barâdarzâde-ye man tâze râh raftan yâd gerefte ast.}
\end{quote}

\begin{quote}
My nephew has recently learned to walk.
\end{quote}

\emph{barâdar} 'brother' + \emph{zâde} 'born of'.

\begin{cousinweb}
\emph{barâdarzâde} is built on two inherited words: \emph{barâdar}, 'brother', which EtymDB traces to Middle Persian (EtymDB), and \emph{zâde}, 'born of'.
\end{cousinweb}

\begin{grammarlens}[title={What you've built}]
One more word for telling, explaining and describing at B1.
\end{grammarlens}

\subsection*{Your turn}
\begin{itemize}
\raggedright
  \item \emph{Say it:} \emph{barâdarzâde}
  \item \emph{Say it:} \emph{barâdarzâde}, once more
  \item \emph{Say it:} a sentence of your own with \emph{barâdarzâde}
  \item \emph{Recall:} say the Persian for a maternal cousin, then the Persian for a sister's child, then say \emph{barâdarzâde} again
\end{itemize}

\begin{tcolorbox}[breakable,colback=teal!4,colframe=teal!35!black,title={Before you move on}]
What does \fa{برادرزاده} mean? (\textquotedblleft{}A brother's child, a nephew or niece\textquotedblright{}.) Say it once more.
\end{tcolorbox}
\section[\texorpdfstring{\fa{مادرشوهر} (mâdarshohar)}{مادرشوهر (mâdarshohar)}]{\fa{مادرشوهر} (mâdarshohar) — a mother-in-law (husband's mother)}

\label{lesson:FA-C501-madarshohar}



\begin{quote}
Before the new one: say the Persian for a sister's child, then the Persian for a brother's child.
\end{quote}

\subsection*{You'll want to know: \fa{مادرشوهر}}
\textbf{\fa{مادرشوهر}} (\emph{mâdarshohar}) — \textquotedblleft{}a mother-in-law (husband's mother)\textquotedblright{}.

\begin{quote}
\fa{مادرشوهرش} \fa{زن} \fa{مهربانی} \fa{است}.
\end{quote}

\begin{quote}
\emph{mâdarshoharash zan-e mehrabâni ast.}
\end{quote}

\begin{quote}
Her mother-in-law is a kind woman.
\end{quote}

\emph{mâdar} 'mother' + \emph{shohar} 'husband'. A wife's mother is \emph{mâdarzan}.

\begin{cousinweb}
\emph{mâdarshohar} is built on two native words: \emph{mâdar}, 'mother', inherited from Middle Persian, with Sanskrit \emph{mātṛ} and Russian \emph{mat'} among its cousins (EtymDB), and \emph{shohar}, 'husband'.
\end{cousinweb}

\begin{grammarlens}[title={What you've built}]
One more word for telling, explaining and describing at B1.
\end{grammarlens}

\subsection*{Your turn}
\begin{itemize}
\raggedright
  \item \emph{Say it:} \emph{mâdarshohar}
  \item \emph{Say it:} \emph{mâdarshohar}, once more
  \item \emph{Say it:} a sentence of your own with \emph{mâdarshohar}
  \item \emph{Recall:} say the Persian for a sister's child, then the Persian for a brother's child, then say \emph{mâdarshohar} again
\end{itemize}

\begin{tcolorbox}[breakable,colback=teal!4,colframe=teal!35!black,title={Before you move on}]
What does \fa{مادرشوهر} mean? (\textquotedblleft{}A mother-in-law (husband's mother)\textquotedblright{}.) Say it once more.
\end{tcolorbox}
\section[\texorpdfstring{\fa{پدرزن} (pedarzan)}{پدرزن (pedarzan)}]{\fa{پدرزن} (pedarzan) — a father-in-law (wife's father)}

\label{lesson:FA-C501-pedarzan}



\begin{quote}
Before the new one: say the Persian for a brother's child, then the Persian for a mother-in-law.
\end{quote}

\subsection*{You'll want to know: \fa{پدرزن}}
\textbf{\fa{پدرزن}} (\emph{pedarzan}) — \textquotedblleft{}a father-in-law (wife's father)\textquotedblright{}.

\begin{quote}
\fa{او} \fa{هر} \fa{هفته} \fa{با} \fa{پدرزنش} \fa{شطرنج} \fa{بازی} \fa{میکند}.
\end{quote}

\begin{quote}
\emph{u har hafte bâ pedarzanash shatranj bâzi mikonad.}
\end{quote}

\begin{quote}
He plays chess with his father-in-law every week.
\end{quote}

\emph{pedar} 'father' + \emph{zan} 'wife'. A husband's father is \emph{pedarshohar}.

\begin{cousinweb}
\emph{pedarzan} is built on two inherited words: \emph{pedar}, 'father', which EtymDB traces to Middle Persian (EtymDB), and \emph{zan}, 'woman, wife'.
\end{cousinweb}

\begin{grammarlens}[title={What you've built}]
One more word for telling, explaining and describing at B1.
\end{grammarlens}

\subsection*{Your turn}
\begin{itemize}
\raggedright
  \item \emph{Say it:} \emph{pedarzan}
  \item \emph{Say it:} \emph{pedarzan}, once more
  \item \emph{Say it:} a sentence of your own with \emph{pedarzan}
  \item \emph{Recall:} say the Persian for a brother's child, then the Persian for a mother-in-law, then say \emph{pedarzan} again
\end{itemize}

\begin{tcolorbox}[breakable,colback=teal!4,colframe=teal!35!black,title={Before you move on}]
What does \fa{پدرزن} mean? (\textquotedblleft{}A father-in-law (wife's father)\textquotedblright{}.) Say it once more.
\end{tcolorbox}
\section[\texorpdfstring{\fa{دوقلو} (doqolu)}{دوقلو (doqolu)}]{\fa{دوقلو} (doqolu) — twins; a twin}

\label{lesson:FA-C501-doqolu}



\begin{quote}
Before the new one: say the Persian for a mother-in-law, then the Persian for a father-in-law.
\end{quote}

\subsection*{You'll want to know: \fa{دوقلو}}
\textbf{\fa{دوقلو}} (\emph{doqolu}) — \textquotedblleft{}twins; a twin\textquotedblright{}.

\begin{quote}
\fa{همسایه} \fa{ما} \fa{دوقلو} \fa{دارد}.
\end{quote}

\begin{quote}
\emph{hamsâye-ye mâ doqolu dârad.}
\end{quote}

\begin{quote}
Our neighbour has twins.
\end{quote}

\emph{do} 'two' + \emph{-qolu}. Triplets are \emph{seqolu}.

\begin{cousinweb}
\emph{doqolu} begins with inherited \emph{do}, 'two', but no source used here explains the second part, and its origin is not settled.
\end{cousinweb}

\begin{grammarlens}[title={What you've built}]
One more word for telling, explaining and describing at B1.
\end{grammarlens}

\subsection*{Your turn}
\begin{itemize}
\raggedright
  \item \emph{Say it:} \emph{doqolu}
  \item \emph{Say it:} \emph{doqolu}, once more
  \item \emph{Say it:} a sentence of your own with \emph{doqolu}
  \item \emph{Recall:} say the Persian for a mother-in-law, then the Persian for a father-in-law, then say \emph{doqolu} again
\end{itemize}

\begin{tcolorbox}[breakable,colback=teal!4,colframe=teal!35!black,title={Before you move on}]
What does \fa{دوقلو} mean? (\textquotedblleft{}Twins; a twin\textquotedblright{}.) Say it once more.
\end{tcolorbox}
\section[\texorpdfstring{\fa{فرزند} (farzand)}{فرزند (farzand)}]{\fa{فرزند} (farzand) — a child (son or daughter), offspring}

\label{lesson:FA-C501-farzand}



\begin{quote}
Before the new one: say the Persian for a father-in-law, then the Persian for twins.
\end{quote}

\subsection*{You'll want to know: \fa{فرزند}}
\textbf{\fa{فرزند}} (\emph{farzand}) — \textquotedblleft{}a child (son or daughter), offspring\textquotedblright{}.

\begin{quote}
\fa{آنها} \fa{سه} \fa{فرزند} \fa{دارند}.
\end{quote}

\begin{quote}
\emph{ânhâ se farzand dârand.}
\end{quote}

\begin{quote}
They have three children.
\end{quote}

More formal than \emph{bachche}. Plural \emph{farzandân}.

\begin{cousinweb}
\emph{farzand} is inherited Persian: EtymDB lists it beside Avestan \emph{frazanti-}, 'offspring, progeny' (EtymDB).
\end{cousinweb}

\begin{grammarlens}[title={What you've built}]
One more word for telling, explaining and describing at B1.
\end{grammarlens}

\subsection*{Your turn}
\begin{itemize}
\raggedright
  \item \emph{Say it:} \emph{farzand}
  \item \emph{Say it:} \emph{farzand}, once more
  \item \emph{Say it:} a sentence of your own with \emph{farzand}
  \item \emph{Recall:} say the Persian for a father-in-law, then the Persian for twins, then say \emph{farzand} again
\end{itemize}

\begin{tcolorbox}[breakable,colback=teal!4,colframe=teal!35!black,title={Before you move on}]
What does \fa{فرزند} mean? (\textquotedblleft{}A child (son or daughter), offspring\textquotedblright{}.) Say it once more.
\end{tcolorbox}
